% !TEX root = main.tex
% Section 1 — Introduction
% Target venue calibration: npj Digital Medicine (IF 15.1) / Journal of Neural Engineering (IF 3.8)
% Word budget: ~2500 words across six sub-sections.

\section{Introduction}
\label{sec:introduction}

\subsection{From Reactive Detection to Preventive Anticipation}
\label{sec:intro-motivation}

Brain--computer interfaces (BCIs) have, for most of their modern history, been \emph{reactive} systems: neural or physiological signals are acquired, features are computed, and a decoder reports a state that has already occurred or a command that the user has already tried to issue. This detection-centric paradigm has delivered two decades of methodological progress and a growing body of clinical translation, but it leaves a structural blind spot --- the minutes and hours before an adverse neurological or cognitive event, during which intervention is most tractable, are precisely the window the system is not designed to see. Drug-resistant focal epilepsy exemplifies the cost of this blind spot: roughly one-third of epilepsy patients continue to have seizures despite optimal pharmacotherapy, and the sudden, unpredictable recurrence of those seizures is the dominant driver of injury, loss of autonomy, and reduced quality of life \cite{WesalEtAl2026,MalikEtAl2026}. Analogous temporal structures --- a silent pre-event period in which the nervous system has already begun to diverge from its stable regime --- recur in sleep-disordered breathing \cite{JamaraniEtAl2026}, migraine onset \cite{TomalinEtAl2025,KapustynskaEtAl2024,JankevičiūtėEtAl2026}, falls in older adults \cite{MaioraEtAl2024,ChanEtAl2023,Al-HammouriEtAl2025}, and vigilance decrements in safety-critical operation \cite{TheivadasPonnan2025,SchwarzEtAl2023,Al-QuraishiEtAl2024}. The unifying claim of this paper is that a second generation of BCIs --- \emph{predictive} systems that forecast such events, and \emph{preventive} systems that close the loop between forecast and intervention --- is now technically feasible, empirically plausible, and clinically under-evaluated.

The shift from detection to forecasting is not merely a relabelling of classification tasks. It requires that models carry, implicitly or explicitly, a temporal prior over how the neural or physiological state evolves toward an event \cite{GunasekaranEtAl2025,ChenEtAl2026b}. Recent work across machine learning and systems neuroscience converges on the view that predictive structure is itself a biomarker: neural correlates of expectancy \cite{KölblEtAl2025}, pre-stimulus rhythms \cite{EguinoaEtAl2026}, and interoception--exteroception integration \cite{BalarKapila2025} all encode information about what is about to happen, not just what has just happened. Forecasting BCIs exploit this structure to deliver lead time; preventive BCIs translate that lead time into action.

\subsection{Background: Prediction, Preictal States, and Ambulatory Sensing}
\label{sec:intro-background}

Three conceptual ingredients underwrite the predictive--preventive paradigm and anchor the remainder of this paper. The first is the \emph{preictal state} --- an operational label for the period immediately before a seizure (or, by extension, any adverse event) during which neural dynamics statistically diverge from the stable interictal baseline. A decade of quantitative work has established that this divergence is detectable from scalp EEG \cite{GuoEtAl2023,Chen2023,ChoiEtAl2022}, intracranial EEG \cite{WangEtAl2026a,ZhouEtAl2026a,QamarEtAl2025}, and, more recently, wearable multimodal streams \cite{NasseriEtAl2021,YamakawaEtAl2020,HasijaEtAl2025}. The operational definition of the preictal window, however, remains a hyperparameter rather than a settled physiological quantity: \citet{KoutsouvelisEtAl2024} show that seemingly technical choices in preictal-period length move subject-specific performance by 20 percentage points, and that the gap between subject-specific and subject-independent evaluation is comparable in magnitude.

The second ingredient is the distinction between \emph{prediction} and \emph{forecasting}. Prediction, as used in much of the deep-learning literature, treats the task as a binary preictal-versus-interictal classification with sensitivity, specificity, false-prediction rate, and AUC as evaluation metrics. Forecasting, in the clinical-epilepsy sense \cite{CostaEtAl2024,FengEtAl2025,FengEtAl2026}, adds a probabilistic time axis: at each moment the system reports a calibrated risk over a future horizon. The two framings converge computationally but diverge clinically --- calibrated probabilistic forecasting is closer to the form that preventive intervention requires, because an alert that triggers a behavioural, pharmacological, or caregiver response must be acceptable, not merely accurate.

The third ingredient is \emph{ambulatory sensing}. The last five years have seen a decisive shift from controlled-laboratory scalp EEG toward wearable, minimally instrumented, and long-duration recordings \cite{NasseriEtAl2021,YamakawaEtAl2020,YedurkarEtAl2025,Paneru2026,AbbasiEtAl2025,KumarEtAl2026}. This shift is what makes preventive BCIs even conceivable: a system that depends on an epilepsy-monitoring-unit (EMU) admission can detect, but cannot prevent. A system that runs on a smartphone-paired wearable, with on-device inference at milliwatt power budgets \cite{CaoEtAl2026,WangEtAl2021b,LiEtAl2025a,XuYu2026}, can alert the patient, recruit a caregiver, or modify the environment in real time.

\subsection{State of the Art: Prediction Quality Meets Cross-Dataset Reality}
\label{sec:intro-sota}

Against this backdrop, the contemporary literature on predictive BCIs exhibits a striking pattern: within-dataset prediction quality is \emph{very high}, but cross-dataset and ambulatory deployment performance is \emph{substantially lower}, and outcome measurement --- whether the prediction actually prevents anything --- is \emph{rarely reported at all}.

On the prediction-quality axis, the Tier 1 evidence is unambiguous. Multi-teacher knowledge distillation architectures achieve 99.54\% AUC with only 0.082M parameters on CHB-MIT \cite{CaoEtAl2026}; feature-fusion ensembles report 99.34\% sensitivity with a false-prediction rate near 0.01/h \cite{AlkhrijahEtAl2025}; long-horizon attention with neural memory sustains 99.46\% sensitivity across multi-minute horizons \cite{PhamTran2026}; and physics-informed graph neural networks anchored in chimera-state detection generalize across patients at 89.3\% sensitivity, decaying gracefully to 68.2\% at a 90-minute horizon \cite{AmiriEtAl2026}. Diverse architectural families --- convolutional--recurrent hybrids \cite{DaoudBayoumi2019,ChoiEtAl2022,AbdelhameedBayoumi2021}, transformers \cite{MahdiBaghdadi2026,ZhengEtAl2025,WangEtAl2026a,ShiLiu2023,GodoyEtAl2022}, graph networks \cite{ZhouEtAl2026a,AmiriEtAl2026,AbbasiEtAl2025}, and diffusion generators of synthetic preictal data \cite{JiangEtAl2024a,XuEtAl2022,PuahEtAl2025} --- converge on essentially the same point: within-dataset ceilings have been reached.

On the cross-dataset axis, the picture is noticeably less favourable. \citet{JangEtAl2026} report single-channel CHB-MIT performance of 97.92\%, falling to 94.93\% on the SNUH cohort --- a clinically meaningful drop that is nevertheless on the optimistic end of the observed distribution. \citet{LiEtAl2025b} document 5--15 percentage-point degradations from CHB-MIT to the MSSM and Siena datasets; \citet{KoutsouvelisEtAl2024} find up to 20-point gaps between subject-specific and subject-independent evaluation; and scoping reviews across the field \cite{WuEtAl2025,TanEtAl2025,RenEtAl2022,WongEtAl2023,YimanWenqi2025} converge on the conclusion that dataset and evaluation-protocol heterogeneity is a dominant source of reported variance. \citet{AmiriEtAl2026}'s horizon decay --- a 21-percentage-point drop from the near-term window to a 90-minute lead --- is representative: the question is not whether prediction degrades with horizon and domain, but by how much.

Representation-learning advances partially address the generalization problem. Self-supervised pretraining \cite{BanvilleEtAl2020,JiangEtAl2024b}, geometry-aware encoders \cite{PaillardEtAl2025,ChenEtAl2026a}, contrastive objectives \cite{GuoEtAl2023,QiEtAl2025,LiaoEtAl2024}, and generative augmentation \cite{JiangEtAl2024a,XuEtAl2022,PuahEtAl2025,CinquettiEtAl2025} all push toward features that transfer across subjects and datasets. Cross-patient studies --- patient-independent models \cite{DissanayakeEtAl2020,JangEtAl2026}, single-seizure adaptation \cite{TariqEtAl2020,JadduEtAl2024}, patient-adaptive transformers \cite{MahdiBaghdadi2026} --- narrow, but do not close, the subject-specific/subject-independent gap. Efficient architectures make the question of deployment tractable: lightweight statistical models \cite{YangEtAl2025,WangEtAl2025b}, edge-enabled geometric networks \cite{AbbasiEtAl2025}, memory-efficient on-device adaptation \cite{LiEtAl2025a,XuYu2026}, and IoT-integrated BCI stacks \cite{Paneru2026,YedurkarEtAl2025,KumarEtAl2026,ThahniyathEtAl2024} demonstrate that inference can plausibly run on the body, not only on a workstation.

Beyond seizure epilepsy, a smaller but methodologically coherent literature extends the predictive paradigm to adjacent adverse events: non-apneic arousal from multichannel polysomnography \cite{JamaraniEtAl2026}, migraine onset from wearable-derived autonomic and sleep metrics \cite{TomalinEtAl2025,KapustynskaEtAl2024,KapustynskaEtAl2025,JankevičiūtėEtAl2026}, fall risk in older adults \cite{MaioraEtAl2024,ChanEtAl2023,Al-HammouriEtAl2025,CavanaghEtAl2024,GuanEtAl2025,LinEtAl2025,ZhangEtAl2024,LiaoEtAl2025b}, and driver drowsiness or vigilance in safety-critical environments \cite{TheivadasPonnan2025,SchwarzEtAl2023,Al-QuraishiEtAl2024,ZhangEtAl2025}. Scoping reviews \cite{HrubýEtAl2026,MouEtAl2026,MehrlatifanMolla2024,VelasquezEtAl2024} confirm that the predictive paradigm is being re-invented, largely independently, in each of these subfields --- with methodological asymmetry heavily favouring seizure epilepsy in terms of Tier 1 evidence density. Integration-oriented reviews \cite{WesalEtAl2026,QamarEtAl2025} are beginning to converge on the same observation that motivates this paper: the field has many predictors and very few closed preventive loops.

\subsection{The Gap: Prediction Is Measured, Prevention Is Not}
\label{sec:intro-gap}

A careful reading of the evidence above exposes a specific, actionable asymmetry. Across 162 papers systematically screened for this review, essentially every Tier 1 study reports prediction-quality metrics --- sensitivity, specificity, false-prediction rate, AUC, lead time --- on offline benchmarks. Almost none report a \emph{preventive outcome}: a directly measured reduction in the rate, severity, or downstream consequence of the event the system is predicting. The predictive--preventive bridge is asserted in discussion sections \cite{NasseriEtAl2021,YedurkarEtAl2025,WesalEtAl2026} but is rarely traversed in methods sections. The dominant study design is computational or algorithmic; experimental human studies that close the loop between alert and intervention are scarce \cite{WesalEtAl2026}.

The consequences of this asymmetry are not merely rhetorical. First, the field has no calibration point for the magnitude of preventive effect that a non-invasive predictive BCI can plausibly deliver. Invasive neuromodulation devices provide one comparator --- NeuroPace RNS trials report 40--70\% reduction in seizure frequency at two years --- but a non-invasive alert-plus-behavioural-modification loop is a categorically different intervention with unknown yield. Second, without outcome measurement, standard clinical translation pathways (preregistered pilots, sham-controlled trials, regulatory submissions) cannot be entered, because the primary outcome that those pathways require has not been defined, let alone measured, in published work. Third, the absence of outcome data starves follow-on cost-effectiveness analyses \cite{VelasquezEtAl2024} of the quantity they most need.

This paper addresses the gap directly. Rather than propose yet another prediction architecture, we treat prediction as a solved-enough problem --- per the Tier 1 consensus above --- and ask whether a standard predictive model, integrated with a standard behavioural and caregiver intervention protocol, produces a measurable reduction in seizure rate in an ambulatory drug-resistant focal epilepsy cohort. The contribution is \emph{measurement}, not model engineering.

\subsection{Hypothesis and Approach}
\label{sec:intro-hypothesis}

We formalize the claim as a prospective single-centre within-subject pre/post pilot. Twenty adults with drug-resistant focal epilepsy, enrolled at a single epilepsy monitoring unit, wear a validated scalp-EEG wearable and carry a smartphone running a pre-trained CNN--LSTM--attention seizure-prediction model of the class reported in \cite{CaoEtAl2026,AlkhrijahEtAl2025,LiEtAl2025b}. A four-week baseline phase runs the model in silent mode to establish each participant's baseline seizure rate, on-device sensitivity/false-prediction-rate profile, and threshold calibration. A twelve-week intervention phase activates the alert loop: above-threshold predictions trigger a smartphone warning with a behavioural-safety checklist and, concurrently, an SMS to a pre-enrolled caregiver. Our primary hypothesis (H1) is that the within-subject median rate ratio of intervention-phase to baseline-phase seizure rates is $\leq 0.75$ --- a $\geq 25\%$ reduction, set below the invasive-device anchor but above the threshold of clinical irrelevance. Secondary hypotheses bound the prospective prediction quality ($\geq 70\%$ of participants achieving sensitivity $\geq 0.80$ and FPR $\leq 0.20$/h) and link preventive efficacy to prediction lead time, sensitivity, and alert adherence. Pre-registered analyses use Wilcoxon signed-rank and paired negative-binomial regression, with a Spearman--Bonferroni mediation protocol for the lead-time mechanism.

\subsection{Contributions and Outline}
\label{sec:intro-contributions}

This paper makes four contributions.

\paragraph{C1 --- Positional framing.} We articulate a reproducible taxonomy of predictive and preventive BCIs that distinguishes detection, prediction, forecasting, and closed-loop intervention as four distinct methodological commitments, and we document the current imbalance across them.

\paragraph{C2 --- Evidence synthesis.} We synthesize 162 predictive-BCI papers across five SOTA themes (deep sequence models, representation learning, cross-patient generalization, efficient deployment, and cross-domain extension), quantifying the within-dataset/cross-dataset gap, the lead-time decay profile, and the methodological monoculture that currently dominates the field.

\paragraph{C3 --- A testable preventive-outcome protocol.} We specify a within-subject pre/post pilot that measures, for the first time in the non-invasive literature, whether a commodity-architecture predictive BCI coupled to a standard behavioural-plus-caregiver intervention loop reduces seizure rate in a drug-resistant focal epilepsy cohort. The protocol is preregistrable, powered for a Cohen's $d_z = 0.6$ effect at $N=20$, and explicitly isolates loop-closure from model novelty.

\paragraph{C4 --- A follow-on research programme.} The OUT scope of the pilot --- cross-domain extension, factorial generalization studies, sham-controlled arms, and mechanism-dissection factorial RCTs --- is mapped onto a sequenced research programme in which each deferred question becomes the pre-specified next study, conditional on the pilot's feasibility signal.

\paragraph{Paper outline.} Section~\ref{sec:relatedwork} situates the contributions against five SOTA themes and an explicit bridge paragraph contrasting our approach. Section~\ref{sec:methods} (not included in this draft) specifies the study protocol, outcome adjudication, and pre-registered statistical analysis. Section~\ref{sec:discussion} addresses risks, limitations, and the translational pathway implied by the pilot's results.
